\subsection{Two quadratic signed cases}
\label{subsec:signed-quadratic-completion}

\begin{proposition}
The fusion rings \texttt{R7\_4\_1212} and \texttt{R7\_4\_1213} have no
categorification over $\mathbb C$.
\end{proposition}

\begin{proof}
We use the sphericalization conditions and
Proposition~\ref{prop:quadratic-signed-reduction}.
Suppose that a categorification exists, and choose the lifts of positive
second Frobenius--Schur indicator in its sphericalization. Write $T$ for
the signed quotient by the central sign object. Its structure constants
satisfy
\[
 |T_{ij}^k|\leq N_{ij}^k,\qquad T_{ij}^k\equiv N_{ij}^k\pmod2,
 \qquad T_{ii}^j\equiv N_{ii}^j\pmod4.
\]
All basis elements are self-dual, so trace cyclicity and dual reflection
make $T_{ij}^k$ fully symmetric. The nontrivial invertible element $g$
acts positively on all five noninvertible basis elements. The spherical
dimension character $d=(1,1,d_2,\ldots,d_6)$ generates
$K=\mathbb Q(\sqrt5)$, all its coordinates are nonzero, and its codegree
$\delta=\sum_i d_i^2$ belongs to
\[
 \{180,360,270\pm90\sqrt5,450\pm90\sqrt5\}.
\]
Every signed formal codegree lies in $K$, by applying the formal-codegree
field theorem to the sphericalization and dividing its codegrees by two.
It is also an algebraic integer, since it is an eigenvalue of the integer
signed Casimir matrix $\sum_i M_i^2$.
The characters $d$ and $\bar d$ are distinct, including when $\delta$ is
rational. The character $(1,-1,0,\ldots,0)$ has codegree $2$.

Consider multiplication by $X_2$ on the $g=1$ summand, in the basis
$(1+g,X_2,\ldots,X_6)$. Its matrix $M$ is self-adjoint for
$G=\operatorname{diag}(2,1,1,1,1,1)$. Set
\[
 p(t)=\det(tI-M),\qquad
 q(t)=\det(tI-M_{\{1,\ldots,5\},\{1,\ldots,5\}}).
\]
Here the first row of $M$ is $(0,1,0,0,0,0)$ and its first column is
$(0,2,0,0,0,0)^{\mathsf t}$. For an eigenvalue $\lambda$, let $w_\lambda$
be the first diagonal entry of its $G$-orthogonal spectral projection.
The resolvent identity gives
\[
 \frac{q(t)}{p(t)}=\sum_\lambda\frac{w_\lambda}{t-\lambda}.
\]
Each common character of codegree $h$ contributes $2/h$ to this weight.
Indeed its eigenvector has coordinates
$(1,\chi(X_2),\ldots,\chi(X_6))$, whose squared $G$-norm is $h$.
Consequently, if $d_2$ is simple, then
$2p'(d_2)-\delta q(d_2)=0$.

There are exactly $97\,200$ and $259\,200$ possible matrices $M$,
respectively, under the displayed coefficient conditions. We enumerate
only these matrices, without imposing the remaining associativity
equations. For each matrix we test for common roots of $p$ with
\[
 2p'-180q,\quad 2p'-360q,\quad
 (2p'-270q)^2-40500q^2,\quad
 (2p'-450q)^2-40500q^2.
\]
The complete exact calculation, supplied by
\texttt{verify\_signed\_pair.py} and \texttt{one\_matrix.cpp}, is
summarized below. The modular tests use the primes $1009$ and $1013$;
all surviving gcds are subsequently computed over $\mathbb Q$.
Since $p$ is monic, a characteristic-zero common factor cannot disappear
upon reduction modulo either prime.
\[
\begin{array}{c|r|r|r|r}
 \text{ring}&\text{matrices}&\text{first prime}&\text{both primes}
 &\text{simple roots surviving}\\ \hline
 \texttt{R7\_4\_1212}&97\,200&996&372&0\\
 \texttt{R7\_4\_1213}&259\,200&1\,888&288&0
\end{array}
\]
More precisely, every common root remaining over $\mathbb Q$ lies in a
repeated factor of $p$. All repeated factors are linear and have
multiplicity two or three. Thus $d_2$ is rational, and the distinct
characters $d$ and $\bar d$ occupy its eigenspace. Their combined weight is
\[
 b=\frac2\delta+\frac2{\bar\delta}
 =\begin{cases}
  1/45,&\delta=180,\\
  1/30,&\delta=270\pm90\sqrt5,\\
  1/90,&\delta=360\text{ or }450\pm90\sqrt5.
 \end{cases}
\]
For multiplicity two we must have $w_{d_2}=b$. For multiplicity three,
the third character has codegree $h=2/(w_{d_2}-b)$, which must be a
positive integer. Exact rational projections exclude every
multiplicity-three branch and every multiplicity-two branch except
four matrices for \texttt{R7\_4\_1213}, all with $d_2=-3$ and
$\delta=180$.

Each of these four matrices has
\[
 \begin{split}
 p(t)&=(t-2)(t+3)^2(t^3-3t^2-5t+3),\\
 q(t)&=(t+3)(t^4-2t^3-12t^2+12t+13).
 \end{split}
\]
The simple eigenvalue $2$ therefore has codegree $2p'(2)/q(2)=10$.
The cubic factor is irreducible. Its three eigenvalues are simple, so
each corresponding character and codegree is defined over its cubic
field. The codegree also lies in $K$ and hence is rational. Galois
conjugacy makes the three codegrees equal to one positive integer $h$.
The reciprocal-codegree identity then reads
\[
 \frac12+\frac2{180}+\frac1{10}+\frac3h=1,
\]
forcing $h=54/7$, a contradiction.

The supplied verifier rereads the source census, constructs the full
matrix domains, reruns the enumeration, and checks every polynomial gcd
and rational projection. It uses no floating-point computation and
passes unchanged with Python optimization enabled.
\end{proof}
